\appendix
\section{A model case with rooks}\label{sec:torres}

In this appendix we study the construction of the predictor in the rook model, which has only
rows and columns. Everything is explicit in this setting, and one can prove that the Bethe formula
with the one-loop correction and the exact zero-mode normalization approximates the logarithm of the count with error $O(1/k)$,
uniformly from positive densities up to the full board. No proof in the rest of the paper
uses this appendix. We use separate notation here. In particular, $p$ is the Bethe square probability,
$\pi=3.14159\ldots$, and $D$ is a probability law.

\subsection{Definitions and theorem statement}

\paragraph{The model.}
Let $k\ge3$. On the $k\times k$ board, the lines are the $k$ rows and the $k$ columns. For
$\lambda>0$ we set
\[
 Z_k(\lambda)=\sum_{q=0}^{k}\binom kq^2q!\,\lambda^q,
\]
which is the sum of the weight $\lambda^q$ over all placements of nonattacking rooks, where $q$
is the number of rooks. To obtain
the coefficient, first choose the $q$ occupied rows and the $q$ occupied columns. Each
bijection between the two sets then places one rook in each chosen row and gives one of the $q!$
possible placements for this choice. Write $\Lambda=\lambda(k-1)$.

\paragraph{The uniform Bethe fixed point.}
The square probability $p\in(0,1/k)$ and the empty-line probability $y\in(0,1)$ satisfy
\begin{equation}\label{eq:torres-punto}
 p(1-p)=\lambda y^2,\qquad y=1-kp.
\end{equation}
This is the general equation $p_c(1-p_c)^{d-1}=\lambda\prod_{L\ni c}y_L$, with $y_L=1-\sum_{c\in L}p_c$
and $d$ the number of lines through each square, specialized to $d=2$ with all variables equal. Set $\theta=1-y=kp$, the Bethe mean number of rooks per row,
$w=p/(1-p)$, and $u=k\theta/(1-p)$.

\paragraph{The three components of the formula.}
At this fixed point we define
\begin{itemize}
\item the tree term
 \[
  \log Z_B=(\log\lambda)\sum_cp-\sum_Ly\log y-\sum_cp\log p+\sum_c(1-p)\log(1-p),
 \]
 where $c$ ranges over the $k^2$ squares and $L$ over the $2k$ lines,
\item the one-loop matrix $M$, acting on the $2k^2$ states $(c,f)$, with
 $f\in\{\text{row},\text{column}\}$. Its entries are $M[(c,f),(c',f')]=w$ if $f'\ne f$, $c'$ lies on the line of
 family $f'$ through $c$, and $c'\ne c$, and $M[(c,f),(c',f')]=0$ otherwise,
\item the zero-mode factor $\rho_0=\mathsf P/P_G$, where $\mathsf P=\mathbb P(N_R=N_C)$ for two
 independent variables $N_R$ and $N_C$ with law $\mathrm{Bin}(k,y)$, representing the auxiliary numbers of empty rows and
 columns in the Bethe model, and $P_G=(2\pi(V_R+V_C))^{-1/2}$, with $V_R=V_C=ky(1-y)$.
\end{itemize}
The formula and its remainder are
\begin{equation}\label{eq:torres-formula}
 F_k(\lambda)=\log Z_B-\tfrac12\log\det(I+M)+\log\rho_0,\qquad R'_k(\lambda)=\log Z_k(\lambda)-F_k(\lambda).
\end{equation}

\begin{teorema}[Rook model]\label{thm:torres}
The remainder $R'_k$ in~\eqref{eq:torres-formula} has the following properties.
\begin{enumerate}[(i)]
\item \emph{Uniform bound.} For each $\Lambda_0>0$ there is a constant $C=C(\Lambda_0)$ such that, for
 every $k\ge3$ and every $\lambda$ with $\Lambda=\lambda(k-1)\ge\Lambda_0$, we have $|R'_k(\lambda)|\le C/k$.
 Moreover, $R'_k(\lambda)$ has a limit $R'_k(\infty)$ as $\lambda\to\infty$, and the same bound holds
 for this limit.
\item \emph{Full board.}
 \[
 \begin{aligned}
  R'_k(\infty)={}&\log k!-k\log k-k(k-1)\log\bigl(1-\tfrac1k\bigr)
  +\tfrac12(k-1)^2\log\bigl(1-(k-1)^{-2}\bigr)\\
  &+\tfrac12\log\frac{2k}{k-1}+(k-1)\log\frac k{k-1}-\tfrac12\log(4\pi k),
 \end{aligned}
 \]
 and $R'_k(\infty)=-\frac1{12k}-\frac1{4k^2}+O(k^{-3})$.
\item \emph{The condition $\Lambda_0>0$ is necessary.} For each fixed $k$, $R'_k(\lambda)=-\frac12\log(4\pi k^2\lambda)+o(1)$
 as $\lambda\to0^+$. In particular, $R'_k$ is unbounded on $\lambda\in(0,\infty)$.
\end{enumerate}
\end{teorema}

The first three terms in (ii) equal $\log k!-\log\operatorname{perm}_B(J_k)$, where $J_k$ is the
$k\times k$ all-ones matrix and $\operatorname{perm}_B$ is the Bethe permanent~\cite{Vontobel}.

\begin{observacion}[Why $k\ge3$ is needed]
For $k=2$ the factors $1-(k-1)^2w^2$ and $(1-w^2)^{(k-1)^2}$ in Lemma~\ref{lem:torres-lazo} coincide, and
$\det(I+M)=(1-w^2)^2(1+w^2)^2$ vanishes like $y^2$ as $y\to0$. Then
$-\frac12\log\det(I+M)\sim-\log y$ is not canceled by the $\frac12\log y$ term in $\log\rho_0$, and $R'_2\to-\infty$
like $-\frac14\log\Lambda$. For $k\ge3$ only the factor $1-(k-1)^2w^2$ vanishes as $y\to0$. It vanishes
like $y$, and the cancellation is exact (Section~\ref{sec:torres-i}).
\end{observacion}

\begin{observacion}[The constant $C(\Lambda_0)$]
By (iii), for each fixed $k$, $\sup\{|R'_k(\lambda)|:\lambda(k-1)\ge\Lambda_0\}\to\infty$ as
$\Lambda_0\to0$, so $C(\Lambda_0)$ necessarily tends to infinity as $\Lambda_0\to0$. The proof of (i) gives an
explicit bound only for $k\ge k_2(\Lambda_0)$, with $k_2$ very large. For $3\le k<k_2$ the bound follows from a
compactness argument.
\end{observacion}

\paragraph{Idea of the proof.}
An exact identity (Lemma~\ref{lem:torres-identidad}) gives
$Z_k=Z_B\,(1-p)^{-k^2}\,\mathsf P\,\mathbb E_D[q!\,u^{-q}]$, where $\mathbb E_D$ denotes an average under the law
$D$ of the number of occupied rows conditional on $N_R=N_C$. The factor
$(1-p)^{-k^2}\approx e^{k\theta+\theta^2/2}$ is not negligible. In the full-board limit we
have $\mathsf P\to1$, and $D$ concentrates on $q=k$. After division by $\lambda^k$, the identity becomes
the all-ones case of the Watanabe--Chertkov formula~\cite{WC}, which expresses
the permanent of a nonnegative matrix as the product of its Bethe permanent, the permanent of
the matrix with entries $\beta_{ij}(1-\beta_{ij})$, and the factor $\prod_{i,j}(1-\beta_{ij})^{-1}$, where
$\beta$ is the doubly stochastic matrix at the Bethe fixed point. For the all-ones matrix,
$\beta_{ij}=1/k$, and these three factors are the limits of $Z_B/\lambda^k$,
$\mathbb E_D[q!\,u^{-q}]$, and $(1-p)^{-k^2}$. The discrete factor $\mathsf P$, which is the
zero mode, cancels exactly with its counterpart in the formula. What remains is $\mathbb E_D[q!\,u^{-q}]$, the average of a
slowly varying function. This average is evaluated by the one-dimensional Laplace method. The necessary
bounds follow from the fact that $D$ is a Poisson binomial law whose parameters come from the zeros of the
Legendre polynomials, and from Bruns's inequality.
